%!TEX root = ./../../thesis-trier.tex

Consider a constrained synchronization problem \constrsync{L}
with some regular constraint language $L \subseteq \Sigma^*$.

\begin{remark}
	\label{constrsync:rem:alphabet}	
	Every constraint language also specifies an alphabet.
	As we show next, up to $\eqACZEROm$-reductions,
	the specific alphabet does not matter and we can always
	chose a smallest possible alphabet over which to consider the constraint
	language.
	For let $\Gamma \subseteq \Sigma$ and $L \subseteq \Gamma^*$ be regular.
	Then $L \subseteq \Sigma^*$
	and \constrsync{L} with~$L$ over $\Sigma$
	and \constrsync{L} with~$L$ over $\Gamma$
	are equivalent for~$\leACZEROm$-reductions:
	for $\mathcal A$ over $\Sigma$
	we can consider~$\mathcal A_{|\Gamma}$
	and, in the other direction for $\mathcal A$ over~$\Gamma$,
	we can add self-loops for all letters in $\Sigma \setminus \Gamma$ 
	to $\mathcal A$ to obtain an automaton over $\Sigma$.
	Both reductions can be implemented by $\AC^0$-circuits
	by either ignoring certain input bits
	or adding the identity matrix as a transition matrix.
\end{remark}

\IN{constraint PDFA}
In this context, we call a PDFA $\mathcal B$ accepting $L$
a \emph{constraint PDFA}.
%As \constrsync{\emptyset} is solvable in constant time, 
%as no automaton has a synchronizing word in $\emptyset$.
%So, we exclude the empty set from consideration
%as a constraint language, and assume $L \ne \emptyset$
%and that $\mathcal B$ is trim.
Since the constraint languages is assumed to be fixed for a given constrained synchronization 
problem, we can assume $\mathcal B$ is trim.
Furthermore, with respect to Remark~\ref{constrsync:rem:alphabet},
%if not said otherwise
we can choose the alphabet $\Sigma$ 
as small as possible.
So, when nothing else is said, we make the following assumptions.

\begin{enumerate}
	\item The constraint PDFA $\mathcal B$ is trim.
	
	\item For a single regular constraint language $L \subseteq \Sigma^*$,
	every letter is used in $L$, i.e.,
	$\Sigma^* a \Sigma^* \cap L \ne \emptyset$
	for each $a \in \Sigma$.
	
	% nicht immer sinnvoll, z.B. reduktion b^* auf ab^*
	%\item When considering two problems
	%$\constrsync{L_1}$ and $\constrsync{L_2}$
	%over alphabets $\Sigma_{1}$ and $\Sigma_2$ respectively
	%that share some letters
	%and we do not distinguish explicitly the different alphabets,
	%then both are over a
	%a common alphabet, i.e., $\Sigma_1 = \Sigma_2$ 
	%(as we can take $\Sigma_1 \cup \Sigma_2$ otherwise).
\end{enumerate}

Here, PDFAs are only used to describe constraint languages. 
%Hence, all other automata and those appearing
%as inputs to our problems are assumed to be complete. 

Due to \cref{constrsync:rem:alphabet}, we do not specify an alphabet in the problem
specification of \constrsync{L}.
%and might choose whatever alphabet  is most convenient.
Changing the alphabet is actually a (very) special case
of a more general way to encode
the input by so-called solid codes, see Remark~\ref{constrsync:rem:alphabet_solid_special_case}.	


Later, in \cref{rem:PDFA-for-constraint}, we justify the extension of our model to 
PDFAs rather than using only the automaton model with a completely defined transition function introduced earlier.